\chapter{Основания \texorpdfstring{$\Mlayer$}{M}}
\label{ch:foundations}

\section{Состояние и дискретность}

\begin{remark}[Постулат 0.1]
Ранний постулат «решётка существует» заменён \cref{thm:discreteness}.
Дискретность $\Lambda$, шаг $\hl$ и тик $\htick$ выводятся из абсолютных условий
A1--A3, A7, A10, A13, A16, геометрии носителя (§1.6) и бит-бюджета (§3.12.6),
а не вводятся отдельной аксиомой клеточного автомата.
\end{remark}

\begin{theorem}[Дискретность и минимальная ячейка $\hv$]
\label{thm:discreteness}
Допустимый локальный закон $g$ не определён на континууме $x\in\mathbb{R}^3$.
Существует счётный FCC-носитель $\Lambda$, минимальная ячейка $\hv$, шаг $\hl$ и тик $\htick$.
Под-решётка «ниже $\hl$'' не является состоянием $\Mlayer$ --- это запрет дробления регистра,
а не отдельный физический предел.
\end{theorem}

\begin{proof}
Доказательство --- цепочка лемм \ref{lem:01a}--\ref{lem:01e}.
\end{proof}

\begin{lemma}[Конечный алфавит на сайте]
\label{lem:01a}
Из A3, A13 и дискретной арифметики $Z_N[i]$ (§3.12) следует: на каждом сайте
алфавит состояний конечен, пространство конфигураций счётно, а эволюция --- перестановка
конфигураций, а не поток в континууме полей.
\end{lemma}

\begin{proof}
A3 требует сохранения $\sum_x |z(x,t)|^2$ при унитарных шагах.
A13 фиксирует обратимый leapfrog на целых $Z_N[i]$ с проекцией collision.
Следовательно каждая компонента спинора принимает значения из конечного множества
$Z_N[i]\subset\mathbb{C}$, а глобальное состояние --- элемент декартова произведения
конечных множеств по сайтам $\Lambda$.
\end{proof}

\begin{lemma}[FCC и равные light-like рёбра]
\label{lem:01b}
Из A1, A2 и выбора носителя (§1.6) следует: каузальный граф имеет равные
light-like рёбра длины $\hl$, $\Lambda$ --- FCC, число ближайших соседей $N_{12}=12$.
\end{lemma}

\begin{proof}
A1 запрещает за один $\htick$ связь длиннее $\hl$ и вторую координационную оболочку.
A2 локализует $g$ на первой оболочке $N(x)$.
Наиболее плотная изотропная упаковка центров при фиксированном $\hl$ --- FCC;
выпуклая 1-tick оболочка --- кубооктаэдр с $12$ рёбрами длины $\hl$.
\end{proof}

\begin{lemma}[Топологическая атомность ячейки]
\label{lem:01c}
Из A10 и A11 следует: заряд $n\in\mathbb{Z}$ живёт на $\partial(\hv)$;
устойчивый солитон не может занимать «пол-$\hv$'' --- ячейка топологически атомна.
\end{lemma}

\begin{proof}
A10: $\oint d\arg z = 2\pi n$, $n\in\mathbb{Z}$.
A11 запрещает устойчивое размазывание дефекта --- полюс локализуется в $\hv$.
Дробление $\hv$ означало бы новый регистр и иной закон $g$, что противоречит минимальности $\hv$.
\end{proof}

\begin{lemma}[Амплитудно-фазовая сетка с дном]
\label{lem:01d}
Из A7, A5 и §3.12.6 следует бит-бюджет $B_{\hv}=2\pi/\ln 2$, кольцо $N_{\mathrm{ring}}$,
$\Delta\phi_{\min}=\tfrac12$, $z_{\min}>0$.
Дробление $\hv$ --- новый регистр, не тот же $g$.
\end{lemma}

\begin{proof}
A7 задаёт планковский потолок плотности; A5 --- недостижимость $|z|=0$.
§3.12.6 выводит $B_{\hv}$ из фазовой ёмкости ячейки и congruence ladder на $\mathbb{Z}_{N_{\mathrm{ring}}}$.
Минимальная амплитуда $z_{\min}=2^{-B_{\mathrm{amp}}}>0$ --- следствие конечного числа бит амплитуды,
а не fitted параметр.
\end{proof}

\begin{lemma}[Landau-ступень и фиксация $\hl$]
\label{lem:01e}
Из A16 и §5.2.1 следует $s_0=\hbar/2$, $p_0=\hbar/(2\hl)$:
насыщение неопределённости на одной $\hv$ фиксирует $\hl$ при данном $\hbar$ (без отдельного ввода $c$).
\end{lemma}

\begin{proof}
Минимальное действие Arg-события $\Delta\phi_{\min}=\tfrac12$ даёт $s_0=\hbar/2$.
Импульсная ступень $p_0=s_0/\hl$ --- единственная лестница переноса на $\hv$ без дополнительных knobs.
\end{proof}

\begin{corollary}[Дискретность как следствие]
\label{cor:discreteness}
«Дискретность'' --- следствие информации, каузальности, топологии и потолка плотности,
а не postulate «решётка exists''.
Шкала $\hl$ и имя $\ell_P$ --- §7.3--7.4; $\sqrt{\hbar G/c^3}$ --- проверка с CODATA, не определение $\hl$.
\end{corollary}

\section{Состояние и эволюция}

\begin{definition}[Постулат 0.3]
На каждом $\PlanckCell$ спинор $z(x,t)\in Z_N[i]$.
Глобальное состояние $\Psi^t=\{z(x,t)\}_{x\in\Lambda}$ эволюционирует
\[
  \Psi^{t+1}=g(\Psi^t),
\]
локально по $\varepsilon$-окрестности FCC ($N_{12}$; гекс --- срез $\{111\}$).
Один тик $\Mlayer$ --- один $\htick$.
Закон $g$ локален и второго порядка (§3.12).
\end{definition}

\begin{corollary}[Континуум на $\Mlayer$ отсутствует]
\label{cor:no-continuum}
Непрерывные поля, дифференциальные уравнения и «плавное'' время --- слой $\Tlayer$
(readout на \emph{фиксированной} решётке), а не предел $\hl\to 0$.
\end{corollary}

\begin{corollary}[Планковский пол амплитуды]
\label{cor:vacuum-floor}
На $\hv$ нет устремления $|z|\to 0$:
\[
  z_{\min}=2^{-B_{\mathrm{amp}}}=1/Q(\mathrm{frac\_bits}),\qquad
  |z(x)|\ge z_{\min}\ \text{на }\Mlayer.
\]
Абсолютный ноль недостижим (A5); $\rho_{\mathrm{field}}>0$ на вакууме;
фазовая сетка с $\Delta\phi_{\min}=\tfrac12$ и $N_{\mathrm{ring}}$ тоже имеет дно.
\end{corollary}

\section{Закон $g$ из абсолютных условий}

\begin{corollary}[Замыкание $g$]
\label{cor:g-closure}
Допустимый локальный $g$ --- замыкание §3.12: leapfrog второго порядка
и projected collision на $Z_N[i]$, а не postulate «правило CA''.
A1--A16 --- система ограничений, из которой собран $g$.
\end{corollary}

\section{Механика и квантование}

\begin{corollary}[Landau-лестница]
\label{cor:landau}
На $\hv$ нет отдельных postulate для $p$, $L$, $F$, $E$.
Одна ступень $s_0$; остальное --- арифметика $(\hl,\htick)$ и kick-ledger §3.12:
$p_0=s_0/\hl$, $E_0=s_0/\htick$, $L_0=p_0\hl=s_0$, $F_0=p_0/\htick=m_{\arg} g_M$.
Континуумные $p,L,F$ и Madelung $j$ --- $\Tlayer$-readout (§4.3).
\end{corollary}

\begin{corollary}[Полное квантование]
\label{cor:full-quantization}
На $\Mlayer$ нет фундаментального поля $A\sin(kx-\omega t)$.
Любое возбуждение --- целочисленная занятость моды $n_k\in\mathbb{Z}$
(фонон, фотон, пакет $n_E E_0$).
«Волна'' и интерференция --- $\Tlayer$-readout суперпозиции многих квантов.
\end{corollary}

Полный реестр абсолютных условий A1--A16 --- \cref{ch:axioms}.
Закон $g$ --- \cref{ch:evolution}.
